\section{Background \label{ch1:background}}
        \todo{Add more information in here}
        A study examining the influence of a seeker head
        on hypersonic flow was conducted by Li \cite{Li2023}. In this study, Li
        investigated the \gls{BSE} and \gls{OPD} 
        at various seeker-cone angles (ranging from $10^{\circ}$ to
        $60^{\circ}$) and \gls{LOS} angles (ranging from $5^{\circ}$
        to $60^{\circ}$) for a Mach 5 flow using \gls{RANS} simulations, without considering nonequilibrium
        effects. Li's work covers Quadrant I in \Cref{tab:ch1AeroOpticalFidelities}. Additionally, an experimental study of
        optical distortions in a ground-test facility at three different Mach
        numbers ($9$, $12$, and $15$) was performed by Pallotta
        \cite{Pallotta2025}, who found that the maximum distortion occurred at
        Mach 9. Pallotta's work covers Quadrant II in \Cref{tab:ch1AeroOpticalFidelities}. Lynch \cite{Lynch2022} performs an
        experimental study of hypersonic turbulent boundary layers but does not
        consider high enthalpy effects, this work cover covers Quadrant II in
        \Cref{tab:ch1AeroOpticalFidelities}. Castillo
        \cite{CastilloGomez2023a}
        performed an \gls{AO} study on a backward-facing step and a
        compression ramp for a Mach 2 flow using a \gls{LES};
        the results from the simulation were compared against experimental
        results obtained in the Tri-Sonic Wind Tunnel at Notre Dame University.
        On this study Castillo calculated the \gls{OPD} and
        showed optical distortions downstream of the reattachment. Castillo's
        work covers Quadrant II in \Cref{tab:ch1AeroOpticalFidelities}\\.

        Some work showing how much the index of refraction is
        affected by aerodynamic effects has been performed by Gupta et.al.
        \cite{Gupta2020} for supersonic flows and showed that there is indeed,
        a change on index of refraction across a shock wave. For subsonic
        flows, it might be a valid assumption to ignore the aerodynamic effects
        on the optical beam; however for supersonic and specially for
        hypersonic flows this assumption has to be revisit. During a hypersonic
        flight the temperatures behind the shock are very high, that will no
        only induce thermal noise on the sensor but will also affect the
        electromagnetic properties that will influence the sensor's
        performance. Hence a higher distortion as well as power loss will be
        expected in a hypersonic flow. Being able to accurately understand this
        is of importance to build accurate sensors that can not only target
        objects moving at these speeds, but can also accurate predict targets
        from an object moving at hypersonic speeds.\\
